\chapter{Análisis estadístico del corpus}
\label{cap:AnexoG}

Este anexo presenta las principales características estadísticas del corpus CodiEsp~\cite{miranda2020}, que es el conjunto de datos empleado para entrenar y evaluar el clasificador CIE-10.

\section{Descripción general}

CodiEsp es un corpus de textos clínicos en español con sus correspondientes códigos CIE-10-ES anotados manualmente. Para la tarea de diagnóstico se emplean los conjuntos indicados en la Tabla~\ref{tab:codiesp-splits}.

\begin{table}[h]
\centering
\caption{Conjuntos de datos del corpus CodiEsp.}
\label{tab:codiesp-splits}
\begin{tabular}{lrr}
\hline
\textbf{Conjunto} & \textbf{Documentos} & \textbf{Proporción} \\
\hline
Entrenamiento & 500 & 50{,}0\,\% \\
Validación & 250 & 25{,}0\,\% \\
Prueba & 250 & 25{,}0\,\% \\
\hline
\textbf{Total} & \textbf{1.000} & 100\,\% \\
\hline
\end{tabular}
\\[4pt]
{\footnotesize Generada con el \emph{notebook} \texttt{codiesp\_\allowbreak{}analisis\_\allowbreak{}estadistico.ipynb} (\texttt{training/\allowbreak{}data\_analysis/}).}
\end{table}

\section{Estadísticas de los textos clínicos}

Los documentos son informes de alta hospitalaria en español. La Tabla~\ref{tab:text-stats} resume las estadísticas de longitud en número de palabras.

\begin{table}[h]
\centering
\caption{Longitud de los textos clínicos por conjunto (en palabras).}
\label{tab:text-stats}
\begin{tabular}{lrrrrrr}
\hline
\textbf{Conjunto} & \textbf{Media} & \textbf{Des. est.} & \textbf{Mín.} & \textbf{Q1} & \textbf{Mediana} & \textbf{Máx.} \\
\hline
Entrenamiento & 349 & 165 & 73 & 225 & 314 & 1.172 \\
Validación & 352 & 168 & 69 & 225 & 328 & 956 \\
Prueba & 353 & 163 & 90 & 236 & 330 & 965 \\
\hline
\end{tabular}
\\[4pt]
{\footnotesize Generada con el \emph{notebook} \texttt{codiesp\_\allowbreak{}analisis\_\allowbreak{}estadistico.ipynb} (\texttt{training/\allowbreak{}data\_analysis/}).}
\end{table}

La longitud media de 349 palabras puede superar el límite de 512 \emph{tokens} de RigoBERTa-Clinical, dado que la tokenización subpalabra genera típicamente más tokens que palabras. Con el tokenizador del modelo, el 53,6\,\% de los documentos de entrenamiento y el 58\,\% de los de validación y de prueba superan ese límite, lo que implica truncar el final del informe. Los experimentos con ventanas más largas (sección~\ref{sec:entrenamiento}) no mejoraron el resultado, por lo que el impacto sobre el F1 parece limitado.

\section{Estadísticas de las etiquetas}

Cada documento tiene asignados múltiples códigos CIE-10-ES. La Tabla~\ref{tab:label-stats} muestra las estadísticas del número de etiquetas por documento.

\begin{table}[h]
\centering
\caption{Número de etiquetas CIE-10 por documento y conjunto.}
\label{tab:label-stats}
\begin{tabular}{lrrrrrr}
\hline
\textbf{Conjunto} & \textbf{Media} & \textbf{Des. est.} & \textbf{Mín.} & \textbf{Q1} & \textbf{Mediana} & \textbf{Máx.} \\
\hline
Entrenamiento & 11{,}28 & 6{,}95 & 1 & 6 & 10 & 40 \\
Validación & 10{,}71 & 6{,}81 & 1 & 5 & 10 & 33 \\
Prueba & 11{,}37 & 6{,}59 & 1 & 6 & 10 & 35 \\
\hline
\end{tabular}
\\[4pt]
{\footnotesize Generada con el \emph{notebook} \texttt{codiesp\_\allowbreak{}analisis\_\allowbreak{}estadistico.ipynb} (\texttt{training/\allowbreak{}data\_analysis/}).}
\end{table}

En total, el corpus contiene 11.158 asignaciones de código. Los códigos únicos presentes son 2.557, de los cuales 1.767 aparecen en el conjunto de entrenamiento y conforman el espacio de clases del clasificador.

\section{Distribución de códigos}

La distribución de frecuencias de los códigos sigue una ley de potencia: un pequeño número de códigos concentra la mayor parte de las asignaciones, mientras que la mayoría de los 2.557 códigos únicos aparece pocas veces. Los 10 códigos más frecuentes del corpus se indican en la Tabla~\ref{tab:top-codes}.

\begin{table}[h]
\centering
\caption{Diez códigos CIE-10 más frecuentes en CodiEsp.}
\label{tab:top-codes}
\begin{tabular}{llrr}
\hline
\textbf{Código} & \textbf{Descripción} & \textbf{Frecuencia} & \textbf{Porcentaje} \\
\hline
R52 & Dolor, no especificado & 219 & 1{,}96\,\% \\
R69 & Enfermedad NEOM & 198 & 1{,}77\,\% \\
R50.9 & Fiebre, no especificada & 191 & 1{,}71\,\% \\
I10 & Hipertensión esencial (primaria) & 161 & 1{,}44\,\% \\
R59.9 & Adenomegalia, no especificada & 136 & 1{,}22\,\% \\
R60.9 & Edema, no especificado & 124 & 1{,}11\,\% \\
R11.10 & Vómitos, no especificados & 92 & 0{,}82\,\% \\
R59.0 & Adenomegalia localizada & 91 & 0{,}82\,\% \\
B99.9 & Enfermedad infecciosa no especificada & 89 & 0{,}80\,\% \\
R58 & Hemorragia, no clasificada bajo otro concepto & 88 & 0{,}79\,\% \\
\hline
\end{tabular}
\\[4pt]
{\footnotesize Generada con el \emph{notebook} \texttt{codiesp\_\allowbreak{}analisis\_\allowbreak{}estadistico.ipynb} (\texttt{training/\allowbreak{}data\_analysis/}).}
\end{table}

Este desequilibrio extremo motivó el uso de \emph{pos\_weight\_cap} durante el entrenamiento para penalizar los errores en códigos poco frecuentes.

\section{Correlación longitud-etiquetas}

El análisis de correlación entre la longitud del texto y el número de etiquetas asignadas es positivo y moderado:

\begin{itemize}
 \item Correlación longitud (caracteres) -- número de etiquetas: $r = 0{,}557$
 \item Correlación número de palabras -- número de etiquetas: $r = 0{,}528$
\end{itemize}

Los textos más extensos corresponden a ingresos más complejos con más diagnósticos asociados. Esta correlación explica por qué los informes más largos son los más penalizados por la truncación a 512 \emph{tokens}. La alternativa de ventana deslizante se evaluó y no compensó el coste, ya que la fragmentación degrada la coherencia clínica del texto (sección~\ref{sec:entrenamiento}).

\section{Solapamiento entre conjuntos}

La Tabla~\ref{tab:jaccard} muestra el solapamiento de códigos entre los tres conjuntos. El bajo coeficiente de Jaccard refleja la cola larga de la distribución (999 códigos aparecen una sola vez en entrenamiento) y el hecho de que 439 de los 1.143 códigos distintos de prueba (38,4\,\%) no aparecen en entrenamiento. Esto introduce un problema implícito de clasificación sin ejemplos previos que limita el techo teórico del F1-micro: solo 2.345 de las 2.842 asignaciones de código de prueba (82,5\,\%) corresponden a códigos vistos al entrenar.

\begin{table}[h]
\centering
\caption{Solapamiento de códigos entre conjuntos (coeficiente de Jaccard).}
\label{tab:jaccard}
\begin{tabular}{lrrrr}
\hline
\textbf{Par} & \textbf{Únicos A} & \textbf{Únicos B} & \textbf{Comunes} & \textbf{Jaccard} \\
\hline
Entrenamiento vs. Prueba & 1.767 & 1.143 & 704 & 0{,}319 \\
Entrenamiento vs. Validación & 1.767 & 1.158 & 731 & 0{,}333 \\
Prueba vs. Validación & 1.143 & 1.158 & 551 & 0{,}315 \\
\hline
\end{tabular}
\\[4pt]
{\footnotesize Generada con el \emph{notebook} \texttt{codiesp\_\allowbreak{}analisis\_\allowbreak{}estadistico.ipynb} (\texttt{training/\allowbreak{}data\_analysis/}).}
\end{table}

\section{Análisis lingüístico del catálogo CIE-10-ES}
\label{sec:analisis-linguistico}

El análisis lingüístico realizado con SpaCy (modelo \texttt{es\_core\_news\_lg}) sobre las 101.246 entradas del catálogo de diagnósticos CIE-10-ES (\emph{notebook} \texttt{cie10\_spacy\_word\_analysis.ipynb}) proporciona una caracterización del vocabulario médico formal empleado en las descripciones oficiales.

\subsection{Estadísticas generales del vocabulario}

Las cifras generales del vocabulario aparecen en la Tabla~\ref{tab:vocab-stats}.

\begin{table}[h]
\centering
\caption{Estadísticas generales del vocabulario del catálogo CIE-10-ES.}
\label{tab:vocab-stats}
\begin{tabular}{lr}
\hline
\textbf{Métrica} & \textbf{Valor} \\
\hline
Lemas únicos (sin \emph{stopwords}) & 10.874 \\
Media de códigos por lema & 68{,}76 \\
Lemas de especificidad máxima (un único código) & 4.823 (44{,}4\,\%) \\
Lemas transversales (aparecen en múltiples códigos) & 6.051 (55{,}6\,\%) \\
\hline
\end{tabular}
\\[4pt]
{\footnotesize Generada con los \emph{notebooks} \texttt{cie10\_\allowbreak{}spacy\_\allowbreak{}word\_\allowbreak{}analysis.ipynb} y \texttt{cie10\_\allowbreak{}analisis\_\allowbreak{}estadistico.ipynb} (\texttt{training/\allowbreak{}data\_analysis/}).}
\end{table}

El hecho de que casi la mitad de los lemas (44,4\,\%) esté asociada a un único código CIE-10 indica que el catálogo contiene una gran proporción de terminología altamente diagnóstica: si ese término aparece en un informe, la asignación del código correspondiente es casi unívoca. El problema es que, como se analiza más adelante, muchos de estos términos de alta especificidad no aparecen en el lenguaje clínico real.

\subsection{Distribución por categoría gramatical}

Las descripciones oficiales de CIE-10-ES presentan un vocabulario dominado por terminología técnica formal. Los sustantivos y adjetivos suman más del 75\,\% del léxico total (Tabla~\ref{tab:pos-dist}).

\begin{table}[h]
\centering
\caption{Distribución del vocabulario CIE-10-ES por categoría gramatical.}
\label{tab:pos-dist}
\begin{tabular}{lrr}
\hline
\textbf{Categoría gramatical} & \textbf{Lemas} & \textbf{Porcentaje} \\
\hline
Sustantivo & 5.145 & 47{,}3\,\% \\
Adjetivo & 3.214 & 29{,}6\,\% \\
Nombre propio & 1.658 & 15{,}2\,\% \\
Verbo & 484 & 4{,}5\,\% \\
Numeral & 139 & 1{,}3\,\% \\
Otros & 234 & 2{,}1\,\% \\
\hline
\end{tabular}
\\[4pt]
{\footnotesize Generada con el \emph{notebook} \texttt{cie10\_\allowbreak{}spacy\_\allowbreak{}word\_\allowbreak{}analysis.ipynb} (\texttt{training/\allowbreak{}data\_analysis/}) y su salida \texttt{cie10\_\allowbreak{}word\_\allowbreak{}analysis.csv}.}
\end{table}

Los nombres propios (15,2\,\%) corresponden principalmente a epónimos y topónimos clínicos (síndrome de Marfan, enfermedad de Crohn, fiebre del Nilo Occidental), que son marcadores diagnósticos muy específicos pero cuya presencia en el texto clínico real depende de que el médico emplee el epónimo exacto y no una descripción funcional equivalente.

\subsection{Términos de alta especificidad y alta dispersión}

El análisis de asociación entre lemas y códigos permite distinguir dos tipos de términos con distinta utilidad para la clasificación automática.

Los \textbf{términos de alta especificidad} (Tabla~\ref{tab:high-spec}) aparecen asociados a un único código CIE-10. Actúan como marcadores casi exclusivos de un diagnóstico concreto y son los más valiosos para el clasificador de diccionario: su presencia en el texto prácticamente determina la asignación. Sin embargo, corresponden en su mayoría a patologías poco frecuentes (como la crioglobulinemia o las infecciones por la cepa O157 de \emph{E. coli}) cuyo nombre técnico rara vez aparece en la redacción habitual de los informes.

\begin{table}[h]
\centering
\caption{Ejemplos de lemas con especificidad máxima (asociados a un único código CIE-10).}
\label{tab:high-spec}
\begin{tabular}{lrr}
\hline
\textbf{Lema} & \textbf{Frecuencia en catálogo} & \textbf{Códigos distintos} \\
\hline
leucemoide & 7 & 1 \\
crioglobulinemia & 6 & 1 \\
o157 & 6 & 1 \\
\hline
\end{tabular}
\\[4pt]
{\footnotesize Generada con el \emph{notebook} \texttt{cie10\_\allowbreak{}spacy\_\allowbreak{}word\_\allowbreak{}analysis.ipynb} (\texttt{training/\allowbreak{}data\_analysis/}) y su salida \texttt{cie10\_\allowbreak{}word\_\allowbreak{}analysis.csv}.}
\end{table}

Los \textbf{términos de alta dispersión} (Tabla~\ref{tab:high-disp}) aparecen en miles de códigos diferentes; funcionan como modificadores estructurales de las descripciones (lateralidad, temporalidad, tipo de lesión) y no aportan capacidad discriminativa por sí solos.

\begin{table}[h]
\centering
\caption{Lemas con mayor dispersión en el catálogo CIE-10-ES.}
\label{tab:high-disp}
\begin{tabular}{lrr}
\hline
\textbf{Lema} & \textbf{Frecuencia total} & \textbf{Códigos distintos} \\
\hline
contacto & 38.306 & 37.687 \\
especificado & 35.482 & 31.322 \\
sucesivo & 22.039 & 22.039 \\
fractura & 36.297 & 20.481 \\
izquierdo & 15.396 & 15.310 \\
\hline
\end{tabular}
\\[4pt]
{\footnotesize Generada con el \emph{notebook} \texttt{cie10\_\allowbreak{}spacy\_\allowbreak{}word\_\allowbreak{}analysis.ipynb} (\texttt{training/\allowbreak{}data\_analysis/}) y su salida \texttt{cie10\_\allowbreak{}word\_\allowbreak{}analysis.csv}.}
\end{table}

\section{Brecha léxica entre catálogo oficial y lenguaje clínico real}
\label{sec:brecha-lexica}

La validación cruzada del vocabulario del catálogo CIE-10-ES con el corpus CodiEsp revela una diferencia estructural entre el lenguaje de las descripciones oficiales y el lenguaje empleado por los médicos en sus informes hospitalarios (Tabla~\ref{tab:vocab-gap}).

\begin{table}[h]
\centering
\caption{Comparativa de vocabulario entre catálogo CIE-10-ES y corpus CodiEsp.}
\label{tab:vocab-gap}
\begin{tabular}{lr}
\hline
\textbf{Métrica} & \textbf{Valor} \\
\hline
Lemas únicos en CIE-10-ES & 10.874 \\
Lemas únicos en CodiEsp & 13.642 \\
Lemas comunes & 4.166 \\
Solapamiento (fracción de CIE-10-ES presente en CodiEsp) & 38{,}3\,\% \\
Lemas exclusivos de CIE-10-ES & 6.708 \\
Lemas exclusivos de CodiEsp & 9.476 \\
\hline
\end{tabular}
\\[4pt]
{\footnotesize Generada con el \emph{notebook} \texttt{cie10\_\allowbreak{}spacy\_\allowbreak{}word\_\allowbreak{}analysis.ipynb} (\texttt{training/\allowbreak{}data\_analysis/}) y su salida \texttt{cie10\_\allowbreak{}word\_\allowbreak{}analysis.csv}.}
\end{table}

Solo el 38,3\,\% de los lemas del catálogo oficial aparece también en los textos clínicos reales. Los 6.708 lemas restantes (nomenclatura estandarizada) no forman parte del lenguaje habitual de los informes. Inversamente, los 9.476 lemas exclusivos de CodiEsp corresponden a abreviaturas, epónimos coloquiales, jerga clínica y variantes morfológicas propias del entorno hospitalario que no figuran en las descripciones oficiales.

La distribución gramatical amplía este contraste: el catálogo CIE-10-ES presenta mayor densidad de sustantivos técnicos y adjetivos específicos, mientras que los informes de CodiEsp muestran mayor presencia de verbos de acción y adverbios de temporalidad, propios de la narración clínica («el paciente ingresa», «se evidencia», «previamente diagnosticado de»).

Esta asimetría léxica tiene consecuencias directas para el diseño del sistema:

\begin{enumerate}
 \item Un clasificador basado exclusivamente en coincidencia léxica con el catálogo (como el clasificador de diccionario) solo puede actuar sobre el 38,3\,\% del vocabulario oficial que efectivamente aparece en los informes, y queda ciego ante los términos de la jerga clínica real que no tienen correspondencia directa en el catálogo.
 \item La expansión explícita de abreviaturas clínicas (\texttt{HTA} → «hipertensión arterial», \texttt{DM2} → «diabetes mellitus tipo 2», \texttt{EPOC} → «enfermedad pulmonar obstructiva crónica») implementada en el clasificador de diccionario actúa como puente parcial entre los dos vocabularios.
 \item El clasificador neuronal basado en RigoBERTa-Clinical supera esta limitación estructural porque su preentrenamiento sobre grandes corpus de texto clínico en español le permite relacionar las expresiones del lenguaje médico real con los diagnósticos correspondientes sin depender de coincidencia léxica exacta con las descripciones oficiales.
\end{enumerate}

\section{Capítulos de la CIE-10-ES}
\label{sec:capitulos-cie10}

Las figuras de rendimiento por capítulo (Figura~\ref{fig:chapters}) identifican cada capítulo de la CIE-10-ES por su número romano. La Tabla~\ref{tab:capitulos-cie10} recoge la equivalencia con su nombre y su intervalo de códigos, junto con los pares documento-código del conjunto de prueba que contiene cada uno y el F1-micro obtenido.

{\small
\begin{longtable}{lP{6.6cm}lrr}
\caption{Capítulos de la CIE-10-ES, intervalo de códigos, pares del conjunto de prueba y F1-micro del modelo servido (umbral 0,2). Los nombres salen de \texttt{ai\_engine/model/cie10\_chapters.json}, los intervalos de \texttt{CHAPTER\_RANGES} en \texttt{ai\_engine/classifier.py} y las cifras de \texttt{model/eval\_test.json} (\texttt{make ai-eval-test}). $^{\dagger}$ Menos de 10 pares: excluido de la Figura~\ref{fig:chapters}.}\label{tab:capitulos-cie10}\\
\hline
\textbf{Cap.} & \textbf{Nombre} & \textbf{Códigos} & \textbf{Pares} & \textbf{F1-micro} \\
\hline
\endfirsthead
\hline
\textbf{Cap.} & \textbf{Nombre} & \textbf{Códigos} & \textbf{Pares} & \textbf{F1-micro} \\
\hline
\endhead
I & Ciertas enfermedades infecciosas y parasitarias & A00--B99 & 148 & 0,42 \\
II & Neoplasias & C00--D49 & 138 & 0,30 \\
III & Enfermedades de la sangre y órganos hematopoyéticos & D50--D89 & 82 & 0,52 \\
IV & Enfermedades endocrinas, nutricionales y metabólicas & E00--E89 & 121 & 0,60 \\
V & Trastornos mentales y del comportamiento & F01--F99 & 51 & 0,38 \\
VI & Enfermedades del sistema nervioso & G00--G99 & 51 & 0,20 \\
VII & Enfermedades del ojo y sus anexos & H00--H59 & 115 & 0,44 \\
VIII & Enfermedades del oído y de la apófisis mastoides & H60--H95 & 1 & 0,00$^{\dagger}$ \\
IX & Enfermedades del sistema circulatorio & I00--I99 & 187 & 0,51 \\
X & Enfermedades del sistema respiratorio & J00--J99 & 67 & 0,51 \\
XI & Enfermedades del aparato digestivo & K00--K95 & 178 & 0,49 \\
XII & Enfermedades de la piel y del tejido subcutáneo & L00--L99 & 59 & 0,50 \\
XIII & Enfermedades del sistema osteomuscular & M00--M99 & 66 & 0,36 \\
XIV & Enfermedades del aparato genitourinario & N00--N99 & 165 & 0,41 \\
XV & Embarazo, parto y puerperio & O00--O9A & 1 & 0,00$^{\dagger}$ \\
XVI & Ciertas afecciones originadas en el período perinatal & P00--P96 & 1 & 0,00$^{\dagger}$ \\
XVII & Malformaciones congénitas & Q00--Q99 & 12 & 0,11 \\
XVIII & Síntomas, signos y hallazgos anormales & R00--R99 & 781 & 0,59 \\
XIX & Traumatismos, envenenamientos y otras consecuencias & S00--T88 & 42 & 0,50 \\
XX & Causas externas de morbilidad y mortalidad & V00--Y99 & 7 & 0,67$^{\dagger}$ \\
XXI & Factores que influyen en el estado de salud & Z00--Z99 & 72 & 0,32 \\
XXII & Códigos para propósitos especiales & U00--U85 & 0 & sin pares \\
\hline
\end{longtable}
}
